In this review, we report recent progress in the application of neural-network quantum states to nuclear many-body problems. Starting from the pioneering studies of the deuteron~\cite{keeble_machine_2020} and very light nuclei~\cite{adams_variational_2021}, the field has advanced enormously over the past five years. Thanks to algorithmic developments and the increased availability of GPU resources, it is now possible to compute ground-state properties of medium-mass nuclei~\cite{gnech_distilling_2024} and infinite nuclear matter, simulating up to $42$ particles in periodic boundary conditions~\cite{fore_investigating_2024}. While simulations of these relatively large systems have so far been carried out using ``essential'' nuclear interactions~\cite{schiavilla_two-_2021}, NQS capable of tackling high-resolution potentials derived within chiral EFT have recently been successfully developed~\cite{wen_neural-network_2025,yang_zemach_2025}. Notably, these applications are not limited to ground-state properties, but also extend to neutron–nucleus scattering~\cite{yang_chiral_2025}.

In addition to extending the reach of conventional continuum QMC methods to medium-mass nuclei, the use of NQS has enabled the study of properties of nuclear systems that are inaccessible to conventional QMC approaches. A notable example is the self-emergence of nuclear clusters in the crust of neutron stars~\cite{fore_investigating_2024}, which lowers the energy per particle and increases the proton fraction. Previous continuum QMC studies in this density regime were unable to capture this phenomenon due to the fermion-sign problem and the reliance on variational ans\"atze that are only suitable for describing the uniform liquid phase. As a second example, the combination of the Lorentz Integral Transform with the VMC--NQS method enables QMC calculations of electroweak response functions~\cite{parnes_nuclear_2026} in the low-energy regime. This region is particularly challenging for conventional GFMC approaches based on imaginary-time current--current correlators, since the inversion of the corresponding Euclidean responses is an ill-posed problem and, in practice, becomes increasingly unreliable at low energy transfer, even when using Maximum Entropy techniques~\cite{lovato_electromagnetic_2016}. Importantly, the low-energy response is of central relevance for neutrino experiments, especially those involving astrophysical neutrinos, such as supernova and diffuse supernova neutrinos~\cite{scholberg_supernova_2012}.

In this review, we have also highlighted the fertile connections with condensed-matter applications. One-dimensional fermionic systems provide an excellent, yet challenging, testing ground for NQS architectures~\cite{keeble_machine_2023} and for the development of new ones, including alternatives to standard MLPs, such as Kolmogorov--Arnold wave functions~\cite{bedaque_machine_2024}. In addition, the first calculations of the dilute neutron-matter equation of state using NQS were enabled by the introduction of periodic coordinates as network inputs, an approach originally developed for periodic bosonic systems~\cite{pescia_neural-network_2022}. These include one- and two-dimensional interacting quantum gases with Gaussian interactions, as well as $^4$He confined in a one-dimensional geometry. Similarly, the highly expressive MPNN backflow transformation, now commonly employed to describe confined and periodic nuclear systems, was originally proposed in the context of the homogeneous electron gas~\cite{pescia_message-passing_2024}. Analogously, the Pfaffian--Jastrow ansatz, before being applied to elucidate the onset of nuclear clustering in the neutron-star crust~\cite{fore_investigating_2024}, proved to be of critical importance for an accurate description of the unitary Fermi gas~\cite{kim_neural-network_2024}.

While not explicitly discussed in this review, recent advances in NQS applications to the nuclear quantum many-body problem include systematic studies of hypernuclei~\cite{donna_hypernuclei_2025,zhang_machine_2025}. In addition to their intrinsic interest, as evidenced by intense experimental campaigns~\cite{hashimoto_spectroscopy_2006}, an accurate determination of nucleon--hyperon and nucleon--nucleon--hyperon interactions is of critical importance for the description of matter in the inner core of neutron stars~\cite{lonardoni_hyperon_2015}. In particular, a quantitative understanding of the repulsive components of the nucleon--nucleon--hyperon interaction is widely regarded as a key ingredient for resolving the so-called hyperon puzzle, namely the apparent tension between the onset of hyperons in dense matter and the existence of neutron stars with masses of about two solar masses~\cite{bombaci_hyperon_2017}.

Based on the rapid progress achieved over the past few years, VMC--NQS approaches are expected to play an increasingly important role in the nuclear many-body community in the coming years.
Potential future applications include the combination of NQS with eigenvector continuation techniques~\cite{frame_eigenvector_2018,duguet_colloquium_2024} to enable systematic studies of nuclear interactions and the reliable quantification of theoretical uncertainties. In this context, foundation models recently developed in condensed-matter physics are also expected to play an important role~\cite{rende_foundation_2025}. Both eigenvector continuation and foundation models offer efficient avenues for exploring families of nuclear Hamiltonians characterized by varying couplings and regulator choices. These approaches are complementary to, and potentially more accurate than, Gaussian-process emulators, which have already been employed in conjunction with NQS to study the sensitivity of hypernuclear ground-state energies to the strength and range of the $\Lambda$--nucleon--nucleon interaction~\cite{donna_hypernuclei_2025}.

Another promising future direction is the calculation of real-time quantum dynamics, which is of critical importance for an \emph{ab initio} description of scattering, fission, and fusion processes. Compared with time-dependent Hartree--Fock approaches~\cite{simenel_heavy-ion_2018}, time-dependent variational Monte Carlo (tVMC)~\cite{carleo_localization_2012} fully incorporates dynamical correlations at both short and long distances. Applications of tVMC to NQS in the continuum are still in their infancy, even in condensed-matter systems~\cite{nys_ab-initio_2024}. Consequently, beyond their direct relevance for the experimental program, employing these techniques to solve the nuclear time-dependent Schr\"odinger equation would also provide a stringent testbed, as nuclear dynamics requires the simultaneous treatment of discrete bound states and continuum degrees of freedom. Achieving this goal would however widen the reach of this promising \textit{ab initio} approach to the realm of nuclear dynamics, potentially addressing open questions in fusion, low-energy transfer and, eventually, fission reactions. 


